\documentclass[11pt]{article}

\usepackage[margin=0.5in]{geometry}
\usepackage{parskip}

\begin{document}
    
\begin{center}
    
\title{\Large}{Specific Aims Outline}

\subsection*{Cell Biology Seminar}

\author{Nur Ceren Kabakci}

\date{\today}

\end{center}



\begin{center}

{\Large \textbf Investigating the Molecular Determinants of
    Antimicrobial Peptide Selectivity and
    Mammalian Cell Death}

\end{center}


\section*{Scientific Premise}
Antimicrobial peptides (AMPs) can preferentially interact
with bacterial membranes, but some peptides also exhibit
cytotoxicity toward mammalian cells.

Understanding the molecular determinants of membrane
selectivity and the cellular consequences of its loss
is important for engineering antimicrobial peptides
with improved biological specificity.

\section*{Central Hypothesis}
Specific sequence and physicochemical features of AMPs
determine their preferential interactions with bacterial
versus mammalian membranes.

Loss of membrane selectivity may initiate membrane damage
and intracellular stress responses that contribute to
mammalian cell death.

\section*{Specific Aims}

\subsection*{Aim 1: Define the membrane-level determinants
of AMP selectivity.}

\textbf{Objective:} Determine how AMP sequence and
physicochemical properties influence bacterial and
mammalian membrane interactions.

\textbf{Expected Outcome:} Identify molecular features
associated with selective antibacterial activity.


\subsection*{Aim 2: Determine the cellular mechanisms
of mammalian cell injury.}

\textbf{Objective:} Investigate whether mammalian cell
death results from direct membrane disruption or
downstream intracellular responses.

\textbf{Expected Outcome:} Identify the cellular processes
responsible for AMP-induced mammalian cell death.


\subsection*{Aim 3: Engineer AMP variants with improved
membrane selectivity.}

\textbf{Objective:} Use mechanistic findings to guide
sequence modifications that preserve antibacterial
activity while reducing mammalian cytotoxicity.

\textbf{Expected Outcome:} Establish sequence-structure-
function relationships that inform rational AMP engineering.

\end{document}